\section{Higher-terminal forest stability at high girth}
\label{sec:path-high-girth-terminal-forests}

The root agent proposed using a breadth-first tree marker for all ports.
The localization argument of Fable then extends from one terminal path to any even terminal set.
A growing-window argument bounds each distance layer, while an entropy Moore bound supplies a large cubic ball.
All logarithms have base two.

\begin{lemma}[A marker with few vertices in each layer]
\label{lem:path-bfs-terminal-marker}
Let $N$ be a normalized $k$-pole with $k\ge2$, and let $T$ be any even subset of its ports.
Choose a breadth-first tree rooted at one port.
Let $P$ be the union of its root-to-port paths, as a vertex set.
Then $P$ is connected, contains every port, and meets each distance layer in at most $k$ vertices.
Moreover, if $J$ is a uniform $T$-join and $\xi$ is a uniform even subgraph, then
\[
 \Pr(J\text{ is a forest})
 \le\Pr(\xi-P\text{ is a forest}).
\]
\end{lemma}

\begin{proof}
Each root-to-port path has at most one vertex in each distance layer.
Their union gives the asserted layer bound and contains all degree-two vertices.
The union of their tree edges is connected, so it contains a $T$-join $F_0$.
The map $\xi\mapsto\xi\mathbin\triangle F_0$ is a bijection onto the affine $T$-join space.
Its restriction outside $P$ agrees with that of $\xi$.
A forest image therefore implies the forest event outside $P$.
\end{proof}

\begin{lemma}[Layer width from a thin connected marker]
\label{lem:path-thin-marker-layer-width}
There is an absolute constant $K$ with the following property.
Let $N$ be connected and subcubic, and let $P$ be a connected marked set containing every vertex of degree below three.
Suppose $P$ meets each distance layer from a fixed root in at most $k$ vertices, where $k\ge1$.
If
\[
 \Pr(\xi_N-P\text{ is a forest})\ge\frac12+\varepsilon,
 \qquad 0<\varepsilon\le\frac12,
\]
then every full distance layer has size at most
\[
 K\left[
 \varepsilon^{-2}\bigl(\log(4/\varepsilon)\bigr)^{11}
 +(k+1)^2\right].
\]
\end{lemma}

\begin{proof}
Put $\delta=\varepsilon/2$, and use $R_\delta$ from
Lemma~\ref{lem:path-polynomial-local-threshold}.
Set
\[
 \Theta=\max\{2R_\delta,\,2^{100}(k+1)^2\}.
\]
We first exclude every set $R\subseteq V(N)\setminus P$ satisfying
\[
 |R|\ge\Theta,
 \qquad
 |\delta_N R|\le\frac{|R|}{4(\log|R|)^2}.
\]
Indeed, let $C_0$ be the component of $N-R$ containing $P$, and put
$R^+=V(N)\setminus V(C_0)$.
This set contains $R$, avoids $P$, and has connected complement.
Its boundary is contained in $\delta_N R$.
All its vertices have degree three in $N$.
Lemma~\ref{lem:path-disconnected-localization} therefore gives a localized graph of rank
\[
 \rho=\frac{|R^+|+|\delta_N R^+|}{2}+1\ge|R|/2\ge R_\delta.
\]
Its marked cycle has at most $\rho/(\log\rho)^2$ vertices, by exactly the monotonicity and padding comparison in Lemma~\ref{lem:path-direct-good-set}.
The forest event outside $P$ implies the forest event on $R^+$.
Localization and the definition of $R_\delta$ bound the latter probability by
$1/2+\delta$, a contradiction.

Write $m_j$ for the number of unmarked vertices in layer $j$, with $m_j=0$ outside the layer range.
For the union $R(i,i')$ of its unmarked layers from $i$ through $i'$, let $m$ be its order.
Edges join the same or consecutive layers.
The boundary satisfies
\[
 |\delta_N R(i,i')|
 \le3m_{i-1}+3m_{i'+1}+3k(i'-i+3).
\]
The last term bounds all incidences to marked vertices in layers $i-1$ through $i'+1$.

Suppose $m_j\ge\Theta$.
Start with this layer and include a neighboring layer whenever its mass exceeds
$m/(48(\log m)^2)$, where $m$ is the current window mass.
This process stops because the graph is finite.
Let $\mu$ be the final mass and put $L=\log\mu\ge100$.
Every step multiplies the current mass by more than $1+1/(48L^2)$.
The number of steps is at most
\[
 \frac{\ln\mu}{\ln(1+1/(48L^2))}
 \le96L^2\ln\mu<67L^3.
\]
Both remaining neighbor masses are at most $\mu/(48L^2)$.
Consequently
\[
 |\delta_N R(i,i')|
 \le\frac{\mu}{8L^2}+3k(67L^3+3)
 \le\frac{\mu}{4L^2}.
\]
For completeness, the last inequality follows from
\[
 24kL^2(67L^3+3)\le2000kL^5\le2^L.
\]
The first comparison holds for $L\ge100$.
For the second, $\mu\ge2^{100}(k+1)^2$ gives $k\le2^{L/2-50}$, and
$2000L^5\le2^{L/2+50}$ for $L\ge100$.
The latter inequality holds at $100$, and its right-to-left ratio increases thereafter.
Thus the stopped window is an excluded set, a contradiction.

Every layer therefore has fewer than $\Theta$ unmarked vertices and at most $k$ marked vertices.
Finally,
\[
 R_\delta\le K_0\delta^{-2}\bigl(\log(2/\delta)\bigr)^{11}.
\]
Absorb this bound and the fixed coefficient $2^{100}$ into an absolute $K$.
\end{proof}

\begin{theorem}[Uniform affine forest bound for high-girth poles]
\label{thm:path-high-girth-terminal-forests}
There is an absolute constant $C$ such that the following holds.
Let $N$ be a normalized $k$-pole of rank $h\ge2$ and girth $g$, with
\[
 2\le k\le g.
\]
For every even subset $T$ of its ports,
\[
 \frac{\#\{T\text{-joins in }N\text{ that are forests}\}}{2^h}
 \le\frac12+C g^{13/2}2^{-g/4}.
\]
In particular, the excess tends to zero as $g\to\infty$, uniformly in the rank, the port count, and the even terminal set.
\end{theorem}

\begin{proof}
It suffices to prove the assertion for all sufficiently large $g$.
An increase in $C$ then covers the remaining finitely many girths.
Corollary~\ref{cor:four-near-moore-shore-rank} gives
\[
 n=|V(N)|\ge2^{g/2-5}.
\]
Put
\[
 s=\left\lfloor\frac g2-\log k-10\right\rfloor.
\]
For large $g$, $s\ge1$ and $2s<g$.
The union of the radius-$s$ balls about the ports has at most
\[
 3k2^s\le3\,2^{g/2-10}<n
\]
vertices.
Choose a vertex $v$ outside that union.
The radius-$s$ ball about $v$ contains no port.
Its breadth-first expansion is cubic and has no repeated vertex before radius $s$, since $2s<g$.
It therefore contains at least
\[
 2^s\ge\frac{2^{g/2}}{2^{11}k}
\]
vertices.

Use the breadth-first root and marker from Lemma~\ref{lem:path-bfs-terminal-marker}.
Distances from this root vary by at most $2s$ across the ball about $v$.
Thus that ball meets at most $2s+1\le g$ root-distance layers.
Some layer has size at least
\begin{equation}\label{eq:path-high-girth-layer-lower}
 \frac{2^{g/2}}{2^{11}kg}
 \ge\frac{2^{g/2}}{2^{11}g^2}.
\end{equation}

Suppose the affine forest proportion is at least $1/2+\varepsilon$.
Translation and Lemma~\ref{lem:path-thin-marker-layer-width} bound every layer by
\[
 K\left[
 \varepsilon^{-2}\bigl(\log(4/\varepsilon)\bigr)^{11}
 +(k+1)^2\right].
\]
Set $\varepsilon=Cg^{13/2}2^{-g/4}$.
For sufficiently large $g$, this lies in $(0,1/2]$ and
$\log(4/\varepsilon)\le g$.
The last layer bound is then at most
\[
 \frac{K}{C^2}\frac{2^{g/2}}{g^2}+K(g+1)^2.
\]
Choose $C$ with $K/C^2\le2^{-13}$.
For sufficiently large $g$, the polynomial second term is also at most
$2^{-13}2^{g/2}/g^2$.
The resulting upper bound is strictly smaller than
\eqref{eq:path-high-girth-layer-lower}, a contradiction.
This proves the asserted uniform estimate.
\end{proof}

The hypothesis $h\ge2$ is essential for this argument.
A cycle with every vertex a port has rank one and girth $k$.
Every nonempty even-terminal affine class in that pole consists entirely of forests.

\begin{corollary}[A quarter-space bound across a nontrivial minimum cut]
\label{cor:path-high-girth-noncycle-minimum-cut}
There is an absolute constant $C_1$ such that the following holds for all sufficiently large girths $g$.
Let $G$ be a connected loopless cubic multigraph of rank $r$ and girth $g$.
Suppose a minimum cyclic cut has order $k$ and neither of its shores is a cycle.
Then
\[
 \frac{c(G)}{2^r}
 \le\frac14+C_1 g^{13/2}2^{-g/4}.
\]
Consequently, any sequence of such graphs with ranks tending to infinity, no cyclic cut of size at most four, and cycle counts above $2^{r-1}$ must eventually have the following property:
its minimum cyclic cut order equals its girth, and every minimum cyclic cut has a shortest cycle as one shore.
\end{corollary}

\begin{proof}
At sufficiently large girth, $G$ has a cyclic cut and its minimum order satisfies $k\le g$.
The shores of a minimum cyclic cut are connected, and its boundary edges have distinct endpoints on each shore.
Write $N_1,N_2$ for these normalized $k$-poles, and $h_1,h_2$ for their ranks.
A rank-one normalized pole is a cycle.
Thus $h_1,h_2\ge2$, and
\[
 r=h_1+h_2+k-1.
\]
The near-Moore rank bound gives $h_i\ge2^{g/2-5}$.
If $k=1$, no cycle crosses the cut; the internal-cycle estimate below proves the assertion directly.
Assume henceforth that $k\ge2$.

A cycle crossing the cut uses a nonempty even subset $S$ of its edges.
Its restriction to each shore is an $S$-join forest.
Some pairs of such forests fail to form one cycle, so their product is an upper bound for the crossing count in this sector.
Apply Theorem~\ref{thm:path-high-girth-terminal-forests} to each shore, with the ambient girth lower bound $g$.
The proof of that theorem applies unchanged with a girth lower bound in place of the exact girth.
Writing $\eta_g=Cg^{13/2}2^{-g/4}$, each sector contributes at most
\[
 (1/2+\eta_g)^2 2^{h_1+h_2}.
\]
There are $2^{k-1}-1$ nonempty even subsets.
After normalization by $2^r$, all crossing cycles therefore contribute at most
$(1/2+\eta_g)^2$.

The cycles internal to the shores contribute at most
\[
 \frac{2^{h_1}+2^{h_2}}{2^r}
 =2^{-h_2-k+1}+2^{-h_1-k+1}.
\]
The rank lower bounds make this $o(\eta_g)$.
Absorb this term and $\eta_g^2$ into an absolute multiple of $g^{13/2}2^{-g/4}$.
This proves the first assertion.

For the final assertion, Corollary~\ref{cor:path-unbounded-girth-reduction} forces the girth to tend to infinity.
The first assertion then excludes any minimum cyclic cut with two noncycle shores.
Every remaining minimum cut has a cycle shore.
Its boundary order equals the length of that cycle, so it is at least the girth.
Since the minimum cyclic cut order is at most the girth, equality holds, and that shore is a shortest cycle.
\end{proof}
